% !TEX root = ../main.tex
この部では、Lean 4 を「数学者のための難しい証明支援系」としてではなく、ソフトウェアエンジニアが仕様を小さく書いて確認するための言語として導入します。
Lean はプログラミング言語であると同時に、証明を書くための言語でもあります。
この二面性は、仕様と実装の関係を考えるうえで重要です。

最初に学ぶものは、型、値、関数です。
これらは多くのエンジニアにとって馴染みのある概念ですが、Lean では型がより中心的な役割を持ちます。
型は単なるメモリ上の表現ではなく、「この値はどのような約束を満たしているか」を表す枠組みとして使われます。
\texttt{Nat}、\texttt{Bool}、\texttt{String}、\texttt{List}、\texttt{Option} のような基本型を使いながら、関数を定義して評価し、簡単な性質を確認します。

次に、命題と証明を扱います。
Lean では、「この関数は常に同じ結果を返す」「この処理の後も不変条件が保たれる」といった主張を \texttt{Prop} として書けます。
さらに、その主張が正しいことを \texttt{theorem} や \texttt{example} として証明します。
最初に使う証明道具は、\texttt{rfl}、\texttt{simp}、\texttt{rw}、\texttt{calc} などの小さなものに絞ります。
ここでの目的は、証明の専門家になることではなく、「仕様を Lean 上で検査可能な形にする」感覚を身につけることです。

\texttt{rfl}、\texttt{simp}、\texttt{rw}、\texttt{calc} は、以後の証明コードで頻繁に現れる小さな証明用キーワードです。
\begin{itemize}
\item \texttt{rfl} は、定義を展開して両辺が同じ形になる等式を閉じます。
\item \texttt{simp} は、Lean が知っている単純化規則で式を整理します。
\item \texttt{rw} は、証明済みの等式を使って現在のゴールを書き換えます。
\item \texttt{calc} は、等式や不等式の変形を段階的に並べます。
\end{itemize}
厳密には、\texttt{rfl}、\texttt{simp}、\texttt{rw} はタクティクであり、\texttt{calc} は構文として説明されることが多いです。
本書では、コードを読むときの目印としてこれらをキーワード扱いにし、Lean のコード表示でも同じ色で強調します。

自然言語の仕様から、小さく独立した性質を切り出して Lean へ移す練習をします。
税込金額の計算、文字列正規化、残高が負にならない操作などの身近な例を使って、関数、命題、証明の関係を確認します。

この部を読み終えるころには、Lean で小さな関数を書き、その関数について「この性質が成り立つ」と主張して、簡単な証明を書けるようになります。
型、関数、命題、証明の関係は、本書全体で共有する基礎になります。

\section*{この部で扱うこと}

\begin{itemize}
\item Lean 4 の基本的な型と値
\item \texttt{def} による関数定義
\item \texttt{\#eval} による実行確認
\item \texttt{Bool} と \texttt{Prop} の違い
\item \texttt{example}、\texttt{theorem}、\texttt{by} による証明
\item \texttt{rfl}、\texttt{simp}、\texttt{rw}、\texttt{calc} の基本
\item \texttt{structure} によるデータモデル
\item 事前条件、事後条件、不変条件の最小表現
\end{itemize}
